\documentclass[11pt,a4paper]{article}
\usepackage{fontspec,amsmath,amssymb,amsthm,microtype}
\setmainfont{lmroman10-regular.otf}[BoldFont=lmroman10-bold.otf,ItalicFont=lmroman10-italic.otf,BoldItalicFont=lmroman10-bolditalic.otf]
\usepackage{xcolor}
\usepackage[most]{tcolorbox}
\usepackage{mathrsfs}
\providecommand{\lt}{<}
\providecommand{\gt}{>}
\setmonofont{DejaVuSansMono.ttf}[BoldFont=DejaVuSansMono-Bold.ttf,ItalicFont=DejaVuSansMono-Oblique.ttf,BoldItalicFont=DejaVuSansMono-BoldOblique.ttf,Scale=0.86]
\newfontfamily\leanfallback{STIXTwoMath-Regular.otf}[Scale=0.86]
\definecolor{leanrule}{gray}{0.62}
\definecolor{leanback}{gray}{0.962}
\definecolor{leantitle}{gray}{0.22}
\definecolor{leanmuted}{gray}{0.38}
\newtcolorbox{leanbox}[1]{enhanced,breakable,colback=leanback,colframe=leanrule,
  boxrule=0.4pt,arc=0.6mm,left=2.4mm,right=2.4mm,top=1.8mm,bottom=1.8mm,
  before skip=0.8em,after skip=1em,fonttitle=\small\bfseries,coltitle=leantitle,
  colbacktitle=leanback,titlerule=0.3pt,title={#1}}
\newlength{\leancharwidth}
\newenvironment{leancode}{\par\smallskip\begingroup\ttfamily\footnotesize\raggedright
  \settowidth{\leancharwidth}{M}\setlength{\parskip}{0pt}\setlength{\parindent}{0pt}
  \frenchspacing}{\par\endgroup\smallskip}
\newcommand{\leanline}[2]{\hangindent=\dimexpr#1\leancharwidth+4\leancharwidth\relax
  \hangafter=1\noindent\hspace*{#1\leancharwidth}#2\par}
\newcommand{\leansource}[1]{\par{\footnotesize\color{leanmuted}#1}\par}
\newcommand{\leanlabel}[1]{\par\smallskip\noindent{\small\bfseries #1}\par}
\usepackage[margin=25mm]{geometry}
\usepackage[unicode,hidelinks]{hyperref}
\hypersetup{pdfauthor={LangevinLean Project}}
\setlength{\emergencystretch}{3em}
\setlength{\parskip}{0.4em}
\setlength{\parindent}{0pt}
\allowdisplaybreaks[2]
\raggedbottom
\title{\texorpdfstring{Convergence of the Verlet–Ornstein–Uhlenbeck splitting}{Convergence of the Verlet–Ornstein–Uhlenbeck splitting}}
\author{LangevinLean Project}
\date{}
\begin{document}\maketitle
Draft mathematical exposition; these notes may contain mistakes.


Let \(\mathcal X\) be a real Banach space, \(\mathcal X\times\mathcal X\) with the maximum norm, \(g:\mathcal X\to\mathcal X\) an \(M\)-Lipschitz map, \(\gamma\ge0\), and \(d(v,x)=(\gamma v+g(x),-v)\) the kinetic drift with unit inverse mass. For a forcing path \(B:\mathbb R\to\mathcal X\) with \(B(0)=0\), a solution of the kinetic integral equation on \([0,T]\) is a continuous \(Y:\mathbb R\to\mathcal X\times\mathcal X\) with


\[Y_t-z_0=-\int_0^td(Y_s)\,ds+(B_t,0)\qquad(0\le t\le T).\]


The splitting chain with step \(h\) is \(X_0=z_0\), \(X_{k+1}=S_h(X_k,N_{kh,h})\), where \(S_h\) is the Verlet–Ornstein–Uhlenbeck step and \(N_{a,h}\) the filtered increment of \(B\) on \([a,a+h]\) (module on the kinetic splitting). This module proves an explicit finite-horizon error bound for \(\|X_n-Y_T\|\) when \(nh=T\), in terms of oscillation moduli of \(t\mapsto d(Y_t)\) and of \(B\), and deduces that \(X_{n+1}\) computed with step \(T/(n+1)\) converges to \(Y_T\) as \(n\to\infty\) whenever \(B\) is continuous.




The argument is pathwise and applies in particular to \(\mathcal X=L^2\) of random vectors with \(B\) the vector Brownian path, where the filtered increments are the Ornstein–Uhlenbeck innovations. It identifies the kinetic diffusion as the limit of a scheme whose steps can be analysed separately (a Verlet step and an exact Ornstein–Uhlenbeck step), which is how the Gibbs invariance of the diffusion is obtained later.


\section{\texorpdfstring{The splitting chain driven by one forcing path}{The splitting chain driven by one forcing path}}


\textbf{Definition.} For \(z_0\in\mathcal X\times\mathcal X\), a path \(B\), a step \(h\) and \(k\in\mathbb N\),


\[\begin{aligned}
X_0&=z_0,\qquad X_{k+1}=S_h\bigl(X_k,N_{kh,h}\bigr),\\
N_{kh,h}&=B(kh+h)-B(kh)\\
&\quad-\gamma\int_0^he^{-\gamma(h-s)}\bigl(B(kh+s)-B(kh)\bigr)\,ds.
\end{aligned}\]


The \(k\)-th step uses the filtered increment of the same path \(B\) on the grid interval \([kh,(k+1)h]\), so the chain and the continuous solution are driven by the same input.

\begin{leanbox}{Lean: \texttt{kineticSplitChain}}
\leansource{\texttt{kineticSplitChain} \quad KineticLangevin/\allowbreak{}KineticSplittingConvergence.lean, lines 12–15}
\begin{leancode}
\leanline{0}{def\ kineticSplitChain\ (g\ :\ E\ →\ E)\ (γ\ h\ :\ ℝ)\ (initial\ :\ E×E)\ (B\ :\ ℝ\ →\ E)\ :\ ℕ\ →\ E×E}
\leanline{2}{|\ 0\ =>\ initial}
\leanline{2}{|\ k+1\ =>\ kineticSplitStep\ g\ γ\ h\ (kineticSplitChain\ g\ γ\ h\ initial\ B\ k)}
\leanline{6}{(ouInputIncrement\ γ\ B\ ((k\ :\ ℝ)*h)\ h)}
\end{leancode}
\leanlabel{How the Lean proof follows the argument}
\texttt{kinetic\allowbreak{}Split\allowbreak{}Chain g \ensuremath{\gamma} h initial B} is defined by recursion on \texttt{\ensuremath{\mathbb N}}: the value at \texttt{0} is \texttt{initial}, and the value at \texttt{k+1} is \texttt{kinetic\allowbreak{}Split\allowbreak{}Step g \ensuremath{\gamma} h (kinetic\allowbreak{}Split\allowbreak{}Chain g \ensuremath{\gamma} h initial B k) (ou\allowbreak{}Input\allowbreak{}Increment \ensuremath{\gamma} B ((k : \ensuremath{\mathbb R})*h) h)}.
\end{leanbox}

\section{\texorpdfstring{A finite-horizon error bound in terms of oscillation moduli}{A finite-horizon error bound in terms of oscillation moduli}}


\textbf{Claim.} Let \(g\) be \(M\)-Lipschitz, \(\gamma\ge0\), \(B(0)=0\), and let \(Y\) be continuous and solve the kinetic integral equation on \([0,T]\). Let \(T\ge0\), \(0\le h\le1\), \(\eta\ge0\), \(R\ge0\) and \(n\in\mathbb N\) with \(nh=T\). Assume \(\|Y_t\|\le R\) for \(t\in[0,T]\) and, for all \(s,t\in[0,T]\) with \(|s-t|\le h\),


\[\|d(Y_s)-d(Y_t)\|\le\eta,\qquad\|B(s)-B(t)\|\le\eta .\]


Then, with \(C=(M+1)(M+2)+\frac{\gamma^2}2+\gamma(M+1)\),


\[\|X_n-Y_T\|\le T\Bigl(hC\bigl(R+\|g(0)\|\bigr)+(1+\gamma)\eta\Bigr)e^{(M+1)^2T}.\]


\textbf{Proof.} Put \(e_k=\|X_k-Y_{kh}\|\), \(q=1+(M+1)^2h\) and \(\varepsilon=hC(R+\|g(0)\|)+(1+\gamma)\eta\ge0\).




\emph{Initial error.} At \(t=0\) the integral equation gives \(Y_0-z_0=(B(0),0)=0\), so \(e_0=\|z_0-Y_0\|=0\).




\emph{One step.} Fix \(k<n\) and put \(a=kh\). Then \(a\) and \(a+h=(k+1)h\) lie in \([0,T]\) because \((k+1)h\le nh=T\), and so does \(a+s\) for \(s\in[0,h]\). Let \(N=N_{a,h}\).



\begin{enumerate}


\item \emph{Exact increment.} Subtracting the integral equations at \(a\) and \(a+h\), and using additivity of the integral over adjacent intervals (the integrand \(s\mapsto d(Y_s)\) is continuous because \(d\) is Lipschitz and \(Y\) is continuous), then the substitution \(s\mapsto a+s\),
\[Y_{a+h}-Y_a=-\int_0^hd(Y_{a+s})\,ds+\bigl(B(a+h)-B(a),0\bigr).\]
With the local defect \(D=\int_0^h\bigl(d(Y_{a+s})-d(Y_a)\bigr)\,ds\) this is the Euler identity \(Y_{a+h}=Y_a-hd(Y_a)+(B(a+h)-B(a),0)-D\).

\item \emph{Local defect.} For \(s\in(0,h]\), \(|a+s-a|=s\le h\), so the integrand of \(D\) has norm at most \(\eta\) and \(\|D\|\le h\eta\).

\item \emph{Filtered increment.} For \(s\in[0,h]\), \(\|B(a+s)-B(a)\|\le\eta\); by the last step of the module on the kinetic splitting, \(\|N-(B(a+h)-B(a))\|\le\gamma h\eta\).

\item \emph{Consistency.} By item 1, \(Y_a-hd(Y_a)+(N,0)-Y_{a+h}=\bigl(N-(B(a+h)-B(a)),0\bigr)+D\), whose norm is at most \(\gamma h\eta+h\eta\) by items 2 and 3 (in the maximum norm \(\|(u,0)\|=\|u\|\)). The local defect estimate of the splitting step, with \(\|Y_a\|\le R\), gives \(\|S_h(Y_a,N)-(Y_a-hd(Y_a)+(N,0))\|\le h^2C(R+\|g(0)\|)\). By the triangle inequality,
\[\|S_h(Y_a,N)-Y_{a+h}\|\le h^2C(R+\|g(0)\|)+(1+\gamma)h\eta=h\varepsilon .\]

\item \emph{Stability.} By the stability of the splitting step, \(\|S_h(X_k,N)-S_h(Y_a,N)\|\le q\,\|X_k-Y_a\|=q\,e_k\).

\end{enumerate}



Since \(X_{k+1}=S_h(X_k,N)\), the triangle inequality and items 4 and 5 give


\[e_{k+1}\le q\,e_k+h\varepsilon\qquad(0\le k<n).\qquad(1)\]


\emph{Summation.} If \(e_0=0\), \(q\ge1\), \(b\ge0\) and \(e_{k+1}\le qe_k+b\) for \(k<n\), then \(e_n\le nbq^n\): by induction, \(e_{k+1}\le q\cdot kbq^k+b\le(k+1)bq^{k+1}\) because \(q^{k+1}\ge1\). Applied to (1), \(e_n\le n\,h\varepsilon\,q^n\). By \(1+u\le e^u\), \(q^n\le\bigl(e^{(M+1)^2h}\bigr)^n=e^{(M+1)^2nh}=e^{(M+1)^2T}\). Since \(nh=T\) and \(Y_{nh}=Y_T\),


\[\|X_n-Y_T\|=e_n\le T\varepsilon\,e^{(M+1)^2T},\]


which is the claim. \(\square\)




The bound separates the deterministic consistency of the scheme (the term \(hC(R+\|g(0)\|)\)) from the regularity of the exact solution and of the forcing (the term \((1+\gamma)\eta\)); the hypothesis \(h\le1\) is needed for the stability and consistency estimates of the Verlet step.

\begin{leanbox}{Lean: \texttt{kineticSplitChain\_\allowbreak{}error\_\allowbreak{}of\_\allowbreak{}modulus}}
\leansource{\texttt{kineticSplitChain\_\allowbreak{}error\_\allowbreak{}of\_\allowbreak{}modulus} \quad KineticLangevin/\allowbreak{}KineticSplittingConvergence.lean, lines 19–128}
\begin{leancode}
\leanline{0}{theorem\ kineticSplitChain\_error\_of\_modulus\ (g\ :\ E\ →\ E)\ (M\ γ\ :\ ℝ≥0)}
\leanline{4}{(hg\ :\ LipschitzWith\ M\ g)\ (initial\ :\ E×E)\ (B\ :\ ℝ\ →\ E)\ (Y\ :\ ℝ\ →\ E×E)}
\leanline{4}{(hB0\ :\ B\ 0=0)\ (hY\ :\ Continuous\ Y)\ (T\ h\ η\ R\ :\ ℝ)}
\leanline{4}{(hT\ :\ 0≤T)\ (hh\ :\ 0≤h)\ (hh1\ :\ h≤1)\ (hη\ :\ 0≤η)\ (hR\ :\ 0≤R)}
\leanline{4}{(n\ :\ ℕ)\ (hpartition\ :\ (n\ :\ ℝ)*h=T)}
\leanline{4}{(heY\ :\ ∀\ t\ ∈\ Icc\ 0\ T,\ Y\ t-initial\ =}
\leanline{6}{-(∫\ s\ in\ 0..t,\ kineticNegativeDrift\ g\ γ\ 1\ (Y\ s))+(B\ t,0))}
\leanline{4}{(hbound\ :\ ∀\ t\ ∈\ Icc\ 0\ T,\ ‖Y\ t‖≤R)}
\leanline{4}{(hmodD\ :\ ∀\ s\ ∈\ Icc\ 0\ T,\ ∀\ t\ ∈\ Icc\ 0\ T,\ dist\ s\ t≤h\ →}
\leanline{6}{‖kineticNegativeDrift\ g\ γ\ 1\ (Y\ s)-kineticNegativeDrift\ g\ γ\ 1\ (Y\ t)‖≤η)}
\leanline{4}{(hmodB\ :\ ∀\ s\ ∈\ Icc\ 0\ T,\ ∀\ t\ ∈\ Icc\ 0\ T,\ dist\ s\ t≤h\ →\ ‖B\ s-B\ t‖≤η)\ :}
\leanline{4}{‖kineticSplitChain\ g\ γ\ h\ initial\ B\ n-Y\ T‖\ ≤}
\leanline{6}{T*(h*(((M\ :\ ℝ)+1)*((M\ :\ ℝ)+2)+(γ\ :\ ℝ)\textasciicircum{}2/2+(γ\ :\ ℝ)*((M\ :\ ℝ)+1))*}
\leanline{8}{(R+‖g\ 0‖)+(1+(γ\ :\ ℝ))*η)*Real.exp\ (((M\ :\ ℝ)+1)\textasciicircum{}2*T)}
\end{leancode}
\leanlabel{How the Lean proof follows the argument}
In \texttt{kinetic\allowbreak{}Split\allowbreak{}Chain\_\allowbreak{}error\_\allowbreak{}of\_\allowbreak{}modulus} the friction \texttt{\ensuremath{\gamma}} and constant \texttt{M} are nonnegative reals (\texttt{\ensuremath{\mathbb R}\ensuremath{\ge}0}), \texttt{hpartition : (n : \ensuremath{\mathbb R})*h=T} is \(nh=T\), \texttt{he\allowbreak{}Y} is the integral equation, \texttt{hbound} the bound by \(R\), and \texttt{hmod\allowbreak{}D}, \texttt{hmod\allowbreak{}B} the two oscillation hypotheses (with \texttt{dist s t \ensuremath{\le} h}).



\begin{enumerate}


\item \texttt{d}, \texttt{C}, \texttt{\ensuremath{\varepsilon}}, \texttt{X} and \texttt{e} are introduced with \texttt{let}; \texttt{hd} is the Lipschitz property of the drift (\texttt{kinetic\allowbreak{}Negative\allowbreak{}Drift\_\allowbreak{}lipschitz}); \texttt{h\allowbreak{}Y0} and \texttt{he0} prove \(Y_0=z_0\) and \(e_0=0\); \texttt{hgrid} shows that grid times lie in \([0,T]\).

\item \texttt{hstep} is (1). Inside it, \texttt{D} is \texttt{ULA.\allowbreak{}local\allowbreak{}Defect d (fun s \ensuremath{\mapsto} Y (a+s)) h}, the local defect; \texttt{h\allowbreak{}D} is item 2, from \texttt{interval\allowbreak{}Integral.\allowbreak{}norm\_\allowbreak{}integral\_\allowbreak{}le\_\allowbreak{}of\_\allowbreak{}norm\_\allowbreak{}le\_\allowbreak{}const}; \texttt{h\allowbreak{}N} is item 3 (\texttt{ou\allowbreak{}Input\allowbreak{}Increment\_\allowbreak{}defect}); \texttt{hinc} is \texttt{ULA.\allowbreak{}langevin\_\allowbreak{}integral\_\allowbreak{}increment} (difference of the integral equation at two times), and \texttt{hinc'} its shifted form via \texttt{interval\allowbreak{}Integral.\allowbreak{}integral\_\allowbreak{}comp\_\allowbreak{}add\_\allowbreak{}left}; \texttt{hid} is the Euler identity \texttt{ULA.\allowbreak{}local\allowbreak{}Defect\_\allowbreak{}euler\_\allowbreak{}identity}; \texttt{heuler} is the identity of item 4, checked componentwise by \texttt{module}; \texttt{h\allowbreak{}S} is the consistency bound, combining \texttt{kinetic\allowbreak{}Split\allowbreak{}Step\_\allowbreak{}euler\_\allowbreak{}defect} with the triangle inequality in a \texttt{calc} block; \texttt{hstab} is item 5 (\texttt{kinetic\allowbreak{}Split\allowbreak{}Step\_\allowbreak{}norm\_\allowbreak{}sub\_\allowbreak{}le}).

\item \texttt{he} is the summation lemma \texttt{ULA.\allowbreak{}finite\_\allowbreak{}euler\_\allowbreak{}recurrence\_\allowbreak{}bound}; \texttt{hpow} bounds \(q^n\) by \texttt{pow\_\allowbreak{}le\_\allowbreak{}pow\_\allowbreak{}left\ensuremath{{}_{0}}}, \texttt{Real.\allowbreak{}add\_\allowbreak{}one\_\allowbreak{}le\_\allowbreak{}exp} and \texttt{Real.\allowbreak{}exp\_\allowbreak{}nat\_\allowbreak{}mul}; the final \texttt{calc} substitutes \texttt{hpartition}.

\end{enumerate}
\end{leanbox}


Earlier results used here: \cite[Section~3.11.1]{LangevinLean2026ULASC}; \cite[Section~3.11.3]{LangevinLean2026ULASC}; \cite[Section~3.16.4]{LangevinLean2026ULASC}; \cite[Section~3.16.5]{LangevinLean2026ULASC}.

\section{\texorpdfstring{Convergence for a continuous forcing path}{Convergence for a continuous forcing path}}


\textbf{Claim.} Let \(g\) be \(M\)-Lipschitz, \(\gamma\ge0\), \(B\) continuous with \(B(0)=0\), \(T\ge0\), and let \(Y\) be continuous and solve the kinetic integral equation on \([0,T]\). Then, with \(h_n=T/(n+1)\),


\[\lim_{n\to\infty}X^{(h_n)}_{n+1}=Y_T,\]


where \(X^{(h)}\) denotes the splitting chain with step \(h\).




\textbf{Proof.} The continuous function \(t\mapsto\|Y_t\|\) is bounded on the compact interval \([0,T]\), say by \(R\ge0\). Put \(C\) as in the previous step and \(D=T\bigl(C(R+\|g(0)\|)+1+\gamma\bigr)e^{(M+1)^2T}\ge0\). Let \(\varepsilon>0\) and \(\eta=\varepsilon/(D+1)>0\); then \(\eta D<\eta(D+1)=\varepsilon\).




The maps \(t\mapsto d(Y_t)\) (continuous, because \(d\) is Lipschitz) and \(B\) are continuous on the compact interval \([0,T]\), hence uniformly continuous there (Heine–Cantor theorem): there are \(\delta_D,\delta_B>0\) such that \(\|d(Y_s)-d(Y_t)\|\le\eta\) whenever \(s,t\in[0,T]\), \(|s-t|\le\delta_D\), and \(\|B(s)-B(t)\|\le\eta\) whenever \(|s-t|\le\delta_B\). Since \(h_n\to0\), there is \(n_0\) such that for \(n\ge n_0\)


\[h_n\le1,\qquad h_n\le\eta,\qquad h_n\le\delta_D,\qquad h_n\le\delta_B .\]


Fix \(n\ge n_0\). The oscillation hypotheses of the previous step hold with step \(h_n\) and tolerance \(\eta\), and \((n+1)h_n=T\). The previous step gives


\[\begin{aligned}
&\|X^{(h_n)}_{n+1}-Y_T\|\\
&\quad\le T\bigl(h_nC(R+\|g(0)\|)+(1+\gamma)\eta\bigr)e^{(M+1)^2T}\\
&\quad\le T\,\eta\bigl(C(R+\|g(0)\|)+1+\gamma\bigr)e^{(M+1)^2T}\\
&\quad=\eta D<\varepsilon.
\end{aligned}\]


using \(h_n\le\eta\) and \(C(R+\|g(0)\|)\ge0\). As \(\varepsilon>0\) was arbitrary, the sequence converges to \(Y_T\). \(\square\)




The theorem does not assume uniqueness of \(Y\): every continuous solution is the limit of the same splitting chain, so it is in particular unique.

\begin{leanbox}{Lean: \texttt{kineticSplitChain\_\allowbreak{}tendsto}}
\leansource{\texttt{kineticSplitChain\_\allowbreak{}tendsto} \quad KineticLangevin/\allowbreak{}KineticSplittingConvergence.lean, lines 132–189}
\begin{leancode}
\leanline{0}{theorem\ kineticSplitChain\_tendsto\ (g\ :\ E\ →\ E)\ (M\ γ\ :\ ℝ≥0)}
\leanline{4}{(hg\ :\ LipschitzWith\ M\ g)\ (initial\ :\ E×E)\ (B\ :\ ℝ\ →\ E)\ (Y\ :\ ℝ\ →\ E×E)}
\leanline{4}{(hB\ :\ Continuous\ B)\ (hB0\ :\ B\ 0=0)\ (hY\ :\ Continuous\ Y)\ (T\ :\ ℝ)\ (hT\ :\ 0≤T)}
\leanline{4}{(heY\ :\ ∀\ t\ ∈\ Icc\ 0\ T,\ Y\ t-initial\ =}
\leanline{6}{-(∫\ s\ in\ 0..t,\ kineticNegativeDrift\ g\ γ\ 1\ (Y\ s))+(B\ t,0))\ :}
\leanline{4}{Tendsto\ (fun\ n\ :\ ℕ\ ↦\ kineticSplitChain\ g\ γ\ (T/((n+1\ :\ ℕ)\ :\ ℝ))\ initial\ B\ (n+1))}
\leanline{6}{atTop\ ({\leanfallback 𝓝}\ (Y\ T))}
\end{leancode}
\leanlabel{How the Lean proof follows the argument}
\texttt{kinetic\allowbreak{}Split\allowbreak{}Chain\_\allowbreak{}tendsto} states \texttt{Tendsto (fun n : \ensuremath{\mathbb N} \ensuremath{\mapsto} kinetic\allowbreak{}Split\allowbreak{}Chain g \ensuremath{\gamma} (T/\allowbreak{}((n+1 : \ensuremath{\mathbb N}) : \ensuremath{\mathbb R})) initial B (n+1)) at\allowbreak{}Top ({\leanfallback 𝓝} (Y T))}. \texttt{R} comes from \texttt{is\allowbreak{}Compact\_\allowbreak{}Icc.\allowbreak{}exists\_\allowbreak{}bound\_\allowbreak{}of\_\allowbreak{}continuous\allowbreak{}On}; \texttt{C} and \texttt{D} are the constants; \texttt{hmesh} is \(h_n\to0\). After \texttt{Metric.\allowbreak{}tendsto\_\allowbreak{}at\allowbreak{}Top.\allowbreak{}mpr} (the \(\varepsilon\)–\(N\) form of convergence) the tolerance \texttt{\ensuremath{\eta}} is chosen and \texttt{hscale}, \texttt{hsmall} give \(\eta D<\varepsilon\). \texttt{huc\allowbreak{}D}, \texttt{huc\allowbreak{}B} are uniform continuity on \texttt{Icc 0 T} (\texttt{Is\allowbreak{}Compact.\allowbreak{}uniform\allowbreak{}Continuous\allowbreak{}On\_\allowbreak{}of\_\allowbreak{}continuous}), and \texttt{Metric.\allowbreak{}uniform\allowbreak{}Continuous\allowbreak{}On\_\allowbreak{}iff\_\allowbreak{}le} extracts the radii \texttt{\ensuremath{\delta}\allowbreak{}D}, \texttt{\ensuremath{\delta}\allowbreak{}B}. \texttt{hnear} combines four eventual bounds on the mesh, each from \texttt{hmesh.\allowbreak{}eventually (eventually\_\allowbreak{}lt\_\allowbreak{}nhds …)}, and \texttt{eventually\_\allowbreak{}at\allowbreak{}Top} turns them into an index \texttt{N}. Then \texttt{he} instantiates the previous theorem with \(n+1\) steps (the partition identity being proved by \texttt{field\_\allowbreak{}simp}), and \texttt{hb}, \texttt{hbound'} carry out the final comparison.
\end{leanbox}

\section{\texorpdfstring{The limit is the constructed kinetic flow}{The limit is the constructed kinetic flow}}


\textbf{Claim.} Let \(g\) be \(M\)-Lipschitz, \(\gamma\ge0\), \(T\ge0\), \(z_0\in\mathcal X\times\mathcal X\), and let \(B\) be a continuous path on \([0,T]\). Let \(\tilde B(t)=B(\min(\max(t,0),T))-B(0)\) be its extension to \(\mathbb R\), constant outside \([0,T]\) and normalised to vanish at \(0\), and let \(Y\) be the solution of the kinetic integral equation with forcing \(\tilde B\) and initial state \(z_0\) given by the Picard construction. Then the splitting chain driven by \(\tilde B\), with step \(T/(n+1)\) and \(n+1\) steps, converges to \(Y_T\).




\textbf{Proof.} \(\tilde B\) is continuous (a composition of continuous maps) and \(\tilde B(0)=0\). The Picard construction provides a continuous \(Y\) with \(Y_0=z_0\) satisfying the integral equation on \([0,T]\) with forcing \(\tilde B\). The previous step applies. \(\square\)

\begin{leanbox}{Lean: \texttt{kineticSplitChain\_\allowbreak{}tendsto\_\allowbreak{}kineticFlow}}
\leansource{\texttt{kineticSplitChain\_\allowbreak{}tendsto\_\allowbreak{}kineticFlow} \quad KineticLangevin/\allowbreak{}KineticSplittingConvergence.lean, lines 193–201}
\begin{leancode}
\leanline{0}{theorem\ kineticSplitChain\_tendsto\_kineticFlow\ (g\ :\ E\ →\ E)\ (M\ γ\ :\ ℝ≥0)}
\leanline{4}{(hg\ :\ LipschitzWith\ M\ g)\ (T\ :\ ℝ)\ (hT\ :\ 0≤T)\ (initial\ :\ E×E)}
\leanline{4}{(B\ :\ VelocityForcing\ (E\ :=\ E)\ T)\ :}
\leanline{4}{Tendsto\ (fun\ n\ :\ ℕ\ ↦\ kineticSplitChain\ g\ γ\ (T/((n+1\ :\ ℕ)\ :\ ℝ))\ initial}
\leanline{6}{(normalizedForcing\ T\ hT\ B)\ (n+1))\ atTop}
\leanline{6}{({\leanfallback 𝓝}\ (kineticFlow\ g\ M\ γ\ 1\ hg\ T\ hT\ initial\ B\ T))}
\end{leancode}
\leanlabel{How the Lean proof follows the argument}
\texttt{kinetic\allowbreak{}Split\allowbreak{}Chain\_\allowbreak{}tendsto\_\allowbreak{}kinetic\allowbreak{}Flow} takes the forcing \texttt{B : Velocity\allowbreak{}Forcing T} (continuous maps from \([0,T]\) to \texttt{E}), extends it by \texttt{normalized\allowbreak{}Forcing T h\allowbreak{}T B} (the extension \(\tilde B\), using \texttt{proj\allowbreak{}Icc}), and identifies the limit as \texttt{kinetic\allowbreak{}Flow g M \ensuremath{\gamma} 1 hg T h\allowbreak{}T initial B T}. The proof takes the three properties of the flow from \texttt{kinetic\allowbreak{}Flow\_\allowbreak{}spec} and the continuity and normalisation of the forcing from \texttt{continuous\_\allowbreak{}normalized\allowbreak{}Forcing} and \texttt{normalized\allowbreak{}Forcing\_\allowbreak{}zero}, and applies \texttt{kinetic\allowbreak{}Split\allowbreak{}Chain\_\allowbreak{}tendsto}.
\end{leanbox}

\begin{thebibliography}{ULA-LSI}
\bibitem[ULA-SC]{LangevinLean2026ULASC}
LangevinLean Project. \emph{Unadjusted Langevin under Strong Convexity}. LL-01, 2026. Draft report, content revision d50c0f64815b. \url{https://odakyildiz.github.io/LangevinLean/docs/pdf/UnadjustedLangevin/FullProof.pdf}.
\end{thebibliography}

\end{document}
